\chapter{Literature Review}
\label{chap:literature}
\setlength{\parskip}{0.8em}

\section{Overview and Scope}

The field of network intrusion detection has generated a vast research corpus through more than three decades of work moving from the statistical anomaly detectors and rule-based expert systems pioneered by Denning and Anderson back in the late 1980s, through the era of supervised machine learning of the 2000s and 2010s, to the cutting edge of deep learning, gradient-boosting ensembles, and hybrid architectures using supervised, unsupervised, and reinforcement-learning components. We review the literature most relevant to this thesis in a structured way below, across five different theme areas: benchmark datasets, traditional machine learning methods, unsupervised and anomaly detection methods, deep learning, and hybrid ensembles, before concluding with a synthesis signalling the research gaps which this thesis attempts to address.
We searched papers in IEEE Xplore, ACM Digital Library, arXiv, ScienceDirect, SpringerLink and finally Google Scholar for the words “network intrusion detection”, “machine learning cybersecurity”, “anomaly detection network traffic”, “deep learning IDS”, “CICIDS dataset”, SHAP cybersecurity” and “explainable IDS”. Papers were selected that were from between 2009 and 2025 and also those which (a) report quantitative performance evaluation on commonly used benchmark datasets (b) report per-class absolute values of metrics and not just aggregate accuracy (c) discuss interpretability of the model and/or deployment in an operational setting. 280 papers were searched through to yield around 21 papers retaining those most relevant, following, and methodologically strict. \section{Network Intrusion Detection Datasets}
\subsection{Historical Development}
The quality and representativeness of training data is the primary factor influencing intrusion detection system performance: no algorithm, however advanced, can overcome the effect of systematic omissions or biases in the data on which it is trained. The field has relied on a series of benchmark datasets, each created in an attempt to fix the perceived shortcomings of its predecessor but each with their own set of limitations. The evolution of datasets follows the same trajectory as the evolution of threat actors, moving from simulated DoS and probe attacks of the late nineties, to web-application and brute-force attacks of the early 2010s, through to IoT-centric and supply-chain attacks today.
\subsubsection{KDD Cup 99 and its Limitations}
KDD Cup 99, derived from the 1998 DARPA network traffic simulation, was for decades the primary benchmark of choice for IDS research. Generated by simulating nine weeks of network traffic on a subset of a US Air Force local area network, it contains approximately 4.9 million connection records labelled normal or as falling into one of the following four categories: Denial of Service (DoS), Probe, Remote-to-Local (R2L), User-to-Root (U2R)~\cite{tavallaee2009nslkdd}. KDD Cup 99 was subsequently shown to have approximately 78\% duplicate records and its attack simulation was not representative of realistic traffic generation. Just like KDD Cup 99 with KDD, models trained on and validated using KDD Cup 99 alone achieve artificially elevated claims of accuracy that do not transfer to live operating environments~\cite{tavallaee2009nslkdd}.
NSL-KDD
Tavallaee et al.~\cite{tavallaee2009nslkdd} proposed NSL-KDD in 2009, removing duplicate records and randomly subsampling from this almost entirely redundant dataset, to deliver only the unique subset with roughly 20% fewer attack records. Each of the training and test KDD datasets has roughly 148{,}517 training records and 22{,}544 test records, making it comparatively easy to evaluate against with algorithmic techniques. Although widely used, it still remains “stuck” in 1998 in terms of attack techniques and represented networks. As a consequence it is less and less useful for evaluating new approaches to real live-attack detection methods—indeed, few approaches are, given KDD’s age.
CICIDS2017
One of the first datasets with a high realism budget was developed at the Canadian Institute for Cybersecurity to mitigate some of the aforementioned issues. Sharafaldin, Lashkari, and Ghorbani~\cite{sharafaldin2018cicids} produced CICIDS2017 using a purpose-built traffic generation framework---the so-called “B-Profile system”---that simulates genuine user AND user profiles (web browsing, email, file transfer, streaming, VoIP) using real attack tools (LOIC, Hulk, GoldenEye, Slowloris, SQLMap, BurpSuite, Patator, Ares, nmap, Metasploit) worked against by security professionals as a target network of 12 workstations, 2 servers, and 1 firewall. Traffic was captured from a sample of users over five days of July 2017 using CICFlowMeter (a flow-based network monitoring tool). On a per-flow basis, CICFlowMeter extracts over 80 statistical features. Repeatability will suffer from this sampling error, and difficulties in being able to train from different data labelling providers at different performance levels; datasets used to train us are replacements for engineers. The dataset contains 14 distinct classes of traffic, containing seven classes of attack—therefore making it the most widely-used dataset in IDS research today. Its primary strengths are high-quality realistic background traffic, modern attacking diversity, and realistic benign users who are not trying to incur attention from researchers or the models—a delinquent of a dataset. Its chief weaknesses are an absence of the encrypted-traffic scenarios—and datasets beyond standard TLS are borrowing from KDD and so unable to handle semantic analysis of malicious decoys. They introduce natural class imbalance, adequately described in full in Chapter~\ref{chap:methodology}; the class imbalance is they key weakness, having public knowledge of balancing already available makes the dataset poor.
UNSW-NB15
Moustafa et al.~\cite{moustafa2015unswnb15} develop UNSW-NB15, at the Australian Centre for Cyber Security, using the IXIA PerfectStorm tool to generate 2.5 million records of synthetic. It spans 49 features, and attacks fall into nine categories across three primary families of attacks: Fuzzers, Analysis, Backdoor, DoS, Exploits, Generic, Reconnaissance, Shellcode, and Worms (presumably from the exploits family). Its network generation technique is more clever and sophisticated than the network-simulation stack bag held together to plug holes in datasets like KDD.
CIDDS-001 and Bot-IoT
Ring et al.~\cite{ring2019survey} introduce CIDDS-001 (Coburg Intrusion Detection Data Set) specifically to supply background traffic with high fidelity; CIDDS-001 differs from most in that besides the traffic, it has additional user, device, and network service metadata; enabling detection based on users having behaviour modelled separately to exploits.
Bot-IoT is developed to treat the IoT infiltration threats that are picking up steam; obviously beyond standard web or network security, Bot-IoT learns models based on Botnet, DDoS, DoS, and foreground techniques, making an effort to explore IoT due to particular modelling constraints and specific behaviours users have for the devices which are not normally obvious.

\subsection{Comparative Dataset Analysis}

Table~\ref{tab:datasets} provides a structured comparison of the principal benchmark datasets discussed above.

\begin{table}[H]
\centering
\caption{Comparison of Principal Network Intrusion Detection Benchmark Datasets}
\label{tab:datasets}
\begin{tabular}{lp{5.5cm}rrr}
\toprule
\textbf{Dataset} & \textbf{Key Characteristics} & \textbf{Records} & \textbf{Classes} & \textbf{Year} \\
\midrule
KDD Cup 99    & First IDS benchmark; DARPA-based; 78\% duplicates; outdated & 4.9M  & 4 & 1999 \\
NSL-KDD       & Cleaned KDD; better balance; still outdated traffic         & 148K  & 4 & 2009 \\
CICIDS2017    & Realistic; 80+ features; diverse modern attacks; widely used & 2.8M & 14 & 2017 \\
CIC-IDS2018   & Extended CICIDS2017; larger scale; additional IoT scenarios & 16M+ & 7 & 2018 \\
UNSW-NB15     & 2.5M records; 9 attack types; modern generation methodology & 2.5M  & 9 & 2015 \\
CIDDS-001     & Corporate network context; realistic background traffic     & 33M   & 4 & 2019 \\
Bot-IoT       & IoT-specific; botnet, DDoS, data theft scenarios            & 73M+  & 4 & 2019 \\
\bottomrule
\end{tabular}
\end{table}

\section{Traditional Machine Learning Methods in Intrusion Detection}

\subsection{Survey of Classical Algorithms}

the most thorough survey based on traditional ML solutions for cybersecurity was conducted by Buczak & Guven ~\cite{buczak2016survey}. The authors of this study covered decision trees, Random Forests, Naive Bayes, k-NN, SVMs and their variants. As stated in this article, there is no single algorithm that performs best in all datasets and with all different metrics used for evaluation. Decision Trees provide highest degree of interpretability of results, however they suffer from being very prone to overfitting if the input data contains too many high dimensional features. This will require some sort of pruning or limiting the depth of the Tree so it will generalize better. Naive Bayes classifiers have low computational costs and perform reasonably good when features are independent of each other. However, if the features do depend on each other (which is typical for network traffic data), the classifier will generally perform poor. For example, packet length statistics in both the forward and backward direction are generally highly correlated. Therefore, a Naive Bayes classifier would likely produce a low accuracy result in such cases. SVMs generally have higher performance in high dimensional feature space, as well as theoretical backing via margin maximization. However, due to the fact that the training complexity of SVMs scales quadratically with the number of Support Vectors, they are not practical for use in large scale real time detection systems dealing with millions of flows. Due to the fact that K-NN is based on finding the closest neighbour, it suffers severely from the "curse of dimensionality". As d grows, the variance of pairwise distance decreases rapidly, and therefore loses discriminant capability.

\subsection{Random Forest as the Baseline Ensemble Method}

Random Forest ~\cite{tavallaee2009nslkdd} has become the de facto gold-standard baseline ensemble method for tabular data classification tasks, including network intrusion detection. Random Forest builds an ensemble of decision trees by training each tree on bootstrapped samples of the original dataset with a subset of random features selected at each node. Random Forest provides excellent variance reduction without sacrificing bias. It also includes a feature selection mechanism to identify important attributes of the network flows using mean Gini impurity. Moreover, Random Forest is embarrassingly parallel since each tree can be trained separately and thus allows nearly linear scalability on multi-core processors. In the NSL-KDD evaluation by Tavallaee et al.~\cite{tavallaee2009nslkdd}, Random Forest obtained detection rates above 98 \% with false positive rates less than 2 \% for most attack classes, making it the reference point against which current works compare.

However, Random Forest does not exploit second order gradient information nor are its predictions computed using structured additive updates --- two aspects addressed by the gradient-boosting family studied in Section 2.3.4.

\subsection{Assessment of ML for Realistic Operational IDS}

Sommer and Paxson ~\cite{sommer2010outside} proposed a fundamental critical assessment of applying machine learning to network intrusion detection in operational environments. Their assessment identified several structural issues that evaluations performed on laboratory datasets with known benchmarks miss:

\begin{enumerate}
\item \emph{semantic gap:} Machine Learning models learn statistical relationships that may correspond to none semantically meaningful security concepts; hence the generated alarms may be incomprehensible for the analysts to confirm.

\item \emph{high base-rate problem:} Even though a false positive rate of 0.1 \% might seem acceptable on a network handling million of packets per day, it could generate tens of thousand of unnecessary alarms and overwhelm analysts ability to manage alarms.

\item \emph{dataset representativeness:} The dataset that were used as benchmarks were captured in a specific network at a specific time; hence their distribution of traffic patterns may be quite dissimilar to those found in the operational deployment environment.

\item \emph{adversary manipulation:} Network Intrusion Detection Systems face adversaries that deliberately watch how the System responds and then adapt their attacks to avoid being detected.
\end{enumerate}

All these problems directly motivated part of our Research focused on Interpretability in order to fill the semantic gap between model prediction and Analyst comprehension through SHAP and LIME Analysis.

\subsection{Recent Results: Gradient Boosting }

XGBoost was introduced by Chen and Guestrin in 2016 ~\cite{xgboost} followed shortly after by LightGBM by Ke et al. in 2017 ~\cite{lightgbm} and CatBoost by Prokhorenkova et al. in 2018 ~\cite{catboost}. They introduced three new technologies into tabular data classification that improve upon classical gradient boosting : (i) regularized objective functions that explicitly punish model complexity and help to prevent overfitting ; (ii) second order Taylor expansion of the Loss Function to make Split Decisions more informative than first order Gradients ; and (iii) histogram-based computation of Gradients instead of performing an exact sorting of features to drastically decrease Memory Bandwidth Requirements.

Al-Lail et al.~\cite{allail2023ml} published a comparative Study comparing multiple Machine Learning Algorithms to detect Malicious Traffic Flow on Standard IDS Benchmarks. Al-Lail et al. showed that Gradient Boosting Methods always perform better than Random Forest, SVMs, Neural Networks and others when using realistic Imbalanced Datasets. Experiments on UNSW-NB15 performed by Al-Lail et al. demonstrated that XGBoost achieved a Macro F1-score value of 0.91 whereas Random Forest had 0.85 and a well-tuned MLP had 0.79, reflecting the relative rankings seen here on CICIDS2017.

Table~\ref{tab:ml_summary} provides a structured overview of the principal ML approaches discussed in this section.

\begin{table}[H]
\centering
\caption{Summary of Machine Learning Approaches Applied to Network Intrusion Detection}
\label{tab:ml_summary}
\begin{tabular}{lp{3.6cm}p{3.6cm}l}
\toprule
\textbf{Algorithm} & \textbf{Strengths} & \textbf{Limitations} & \textbf{Key Ref.} \\
\midrule
Logistic Regression & Interpretable; fast; probabilistic output & Linear boundary only & \cite{buczak2016survey} \\
Naive Bayes         & Computationally efficient; handles streaming & Independence assumption & \cite{buczak2016survey} \\
Decision Tree       & Fully explainable; handles mixed types & Prone to overfitting & \cite{buczak2016survey} \\
KNN                 & Non-parametric; flexible boundaries & Slow inference; no training & \cite{buczak2016survey} \\
SVM                 & Strong in high-dim spaces; well-grounded & Quadratic training cost & \cite{buczak2016survey} \\
Random Forest       & High accuracy; feature importance; robust & High memory; opaque & \cite{tavallaee2009nslkdd} \\
XGBoost             & Regularised; fast; handles imbalance well & Many hyperparameters & \cite{allail2023ml} \\
LightGBM            & Fastest GB; leaf-wise; GOSS~+~EFB & Less stable on small data & \cite{nordin2024aireview} \\
DNN / MLP           & Complex feature learning; flexible & Black-box; data-hungry & \cite{javaid2016deep} \\
LSTM / RNN          & Temporal sequence modelling & Slow; hyperparameter-sensitive & \cite{yin2017rnn} \\
Autoencoder         & Unsupervised; anomaly score output & High FPR; threshold-sensitive & \cite{shone2018deep} \\
Isolation Forest    & Unsupervised; fast; no label requirement & Approximate; not multiclass & \cite{khraisat2019survey} \\
\bottomrule
\end{tabular}
\end{table}

\section{Anomaly Detection and Unsupervised Methods}

\subsection{Rationale and Principles}
The fundamental weakness of supervised classification in IDS is that they cannot detect attacks they have never seen in their training set: zero-day attacks. If an adversary develops a novel attack technique that is not present in any dataset, then any purely supervised classifier will be unable to detect that pattern, however well it does on the known attacks. This limitation of course is architectural not practical: it follows from the inductive principle on which supervised learning is built, that what the algorithm tries to learn is a better approximation of the conditional distribution it sees in its training data (and anything that is too far from that will be, by construction, undetectable).
To bypass this limitation, anomaly based detection approaches flag deviations from an established model of normal behaviour, without needing access to labelled attack examples. These may tolerate higher false-positive rates in exchange for the ability to flag things that the best possible supervised classifier would not be able to identify.
Anomaly based IDS fall into three paradigms: statistical anomaly detection which models normal behaviour as some parametric statistical distribution and flags observations under that model that would be improbably unlikely; knowledge based anomaly detection which encodes normal behaviour in an explicit set of rules or specifications derived from human domain knowledge; and machine learning based anomaly detection which will learn a model of normal behaviour from data, without explicit rule specification~\cite{khraisat2019survey}. \subsection{Clustering}

Kim, Lee and Kim~\cite{kim2014hybrid} use k-means clustering for network traffic anomalies (sc-)detecting. Demonstrating that a sufficient features combination can provide a detection separation through unsupervised cluster assignment, k-means partitions the space of feature vectors into a measure of k clusters, minimising the within-cluster sum of the squared distances. These are then flagged potential anomalies if their cluster is small, or if they are some large distance from the cluster centroid. Simple, intuitive and very fast tracing, the algorithms performance is nevertheless sensitive to number of clusters (k), the initialisation of the centres of those clusters, and the scale of the features. These all need substantial tuning in the IDS context.

\subsection{Density Based and Statistical}

Zhang and Zulkernine~\cite{zhang2006anomaly} show multidimensional outlier detection (statistical) for anomaly-based network IDS systems, this depends on model parameters fitting normal traffic, as even small multimodal characteristics lead to traditional distance based anomaly detection failing in naive implementations. Khraisat et al.~\cite{khraisat2019survey} survey a number of various suitable anomaly detection framework evaluations of clustering, statistical and neural network approaches against a number of criteria including the detection rate, false-positive rate, speed of the detection algorithm, and finally how well they generalise.

\subsection{Isolation Forest}

Relatively recent is the Isolation Forest algorithm~\cite{khraisat2019survey} which has rather unexpectedly found itself in growing interested and used as an unsupervised anomaly detector for high dimensional tabular data, in a domain of dimensionality sufficient, but not a high feature count. Unlike distance or densite based approaches however will highlight anomalous instances based on how readily they can be parsed, or cast out, these will statistically exist in lesser number, and be easily sequestered by fewer randomly generated binary tests than typical normal instances, hence isolating in less and more binary slices respectively. The actual $s(x,n)$ score is computed from $E[h(x)]$, whose averaging is across an ensemble of random isolation trees, normalised by the expected $E[h(x)]$:
\begin{equation}
s(x, n) = 2^{-E[h(x)]/c(n)}\,; \, c(n) = 2H(n-1) - 2(n-1)/n
\end{equation}
where `$E[h(x)]$` is an ailment count of steps taken to isolate, `H` is the i$^{th}$ harmonic number, and to recap is path length in a random isolation tree, so approximating the a number of cuts needed for isolation of said x points ($n$), such that closer to 1 is a stronger anomaly, and 0.5 represents typical instances. This mechanism makes it very attractive ($O(n \log n)$, `O(log n)` for per point predictions), naturally parameter light and expedient, (just how many trees? and how much contamination?) and favourable to high dimensional unsampled feature spaces without dense estimation, principal component failure mode, related to kernel space degeneracy is sensitivity to irrelevant features which can start to begin to dilute the very restoring or terminating signal. Somewhat removed by the feature selection pipeline described in Chapter~\ref{chap:methodology}.

\section{Deep Learning Cybersecurity}

\subsection{Deep Neural Networks for IDS}

The ability of deep neural networks to automatically extract hierarchical high level features from raw (< 1 process) or minimized classes of preprocessed traffic make them an irresistably seemily direct applicated to tasks that demand network IDS patience, where background knowledge of the data evading signature rules means feature irrelevant parameter construction cannot be done anyway. Knowing there are not bytes packs to discriminate simple stuff Javaid et al.~\cite{javaid2016deep} represent some of the must closest enthusiasts using and apply deep learning, and show that a 5 layer deep autoencoder multi feat and 2/3s time architectures trained in a unsupervised pre stage before sticking labeled examples in, can outpredict vary well penchant higher than standard KDD ML algorithm training on NSL-KDD.

Figure~\ref{fig:mlp_arch} illustrates the generalised MLP architecture used as one of the evaluated models in this research.

\begin{figure}[H]
\centering
\begin{tikzpicture}[
    node distance=1.4cm and 0.9cm,
    layer/.style={draw=NavyBlue, fill=LightBlue, rounded corners=4pt,
                  minimum width=2.0cm, minimum height=0.8cm, font=\small, align=center, text width=1.9cm},
    arr/.style={-Latex, thick, NavyBlue!70}
]
\node[layer, fill=AccentGreen!20, draw=AccentGreen] (inp) {Input\\(74 feat.)};
\node[layer, right=of inp] (h1) {Hidden 1\\(256)};
\node[layer, right=of h1] (h2) {Hidden 2\\(128)};
\node[layer, right=of h2] (h3) {Hidden 3\\(64)};
\node[layer, right=of h3, fill=AccentRed!15, draw=AccentRed] (out) {Output\\(14 cls)};
\foreach \i/\j in {inp/h1,h1/h2,h2/h3,h3/out} \draw[arr] (\i.east)--(\j.west);
\node[below=0.25cm of h1, font=\footnotesize\itshape, color=gray] {ReLU + Dropout};
\node[below=0.25cm of h2, font=\footnotesize\itshape, color=gray] {ReLU + Dropout};
\node[below=0.25cm of h3, font=\footnotesize\itshape, color=gray] {BN + ReLU};
\node[below=0.25cm of out, font=\footnotesize\itshape, color=gray] {Softmax};
\end{tikzpicture}
\caption{Generalised MLP architecture for multiclass network attack classification. The input layer ingests 74 selected flow-level features; three hidden layers (256~/~128~/~64 units) with ReLU activations and dropout regularisation feed a softmax layer that emits a probability distribution over the 14 traffic classes (benign + 13 attack categories).}
\label{fig:mlp_arch}
\end{figure}

\subsection{Autoencoders for Anomaly Detection}

Shone et al.~\cite{shone2018deep}, developed a non-symmetrical deep autoencoder (NDAE), combined with random forests as a classifier, for network intrusion detection. Autoencoders learn unsupervised representations of network flows; if there is no way to reconstruct a flow based on its latent representation, then the flow is labeled as anomalous. Combining an unsupervised method for learning representations with a supervised method for classifying those representations into anomalous/normal produces good results for detecting both known and unknown variants of attacks that have never appeared before during training.

\subsection{Temporal Models with Recurrent Neural Networks.}

Yin et al.~\cite{yin2017rnn}, show how recurrent neural networks (specifically long short-term memory (LSTM)), could be applied to network intrusion detection. Network attacks are typically represented by some type of sequential pattern ---a series of individual flows that each look like they are legitimate, but when taken together indicate malicious behavior. For instance, a denial-of-service (DDoS) attack is characterized by a prolonged high volume of packets; port scanning involves sending probes quickly in succession; and communication from a botnet command server to a botnet client has specific timing characteristics. RNN's such as LSTM's can identify sequential relationships among the individual elements of a flow by retaining or discarding relevant information from previous time steps through their gating functions. Yin showed that in the case of sequential dependent attack classes, LSTM models performed better than static models, however, the models were much slower to train and had more parameters that needed to be tuned.

Mohammadpour et al.~\cite{mohammadpour2022cnn}, also discussed CNN's, showing that one dimensional CNN's applied to feature vectors can find local correlation between features much faster than fully connected architectures and therefore offer an effective alternative between MLP's and LSTM's.

\subsection{Shortcomings of Applying Deep Learning to IDS. }

Although deep learning techniques in IDS provide significant improvements in terms of detection rate, there remain several limitations to their application:
\begin{itemize}
\item \textbf{Training Data Requirements. }Deep networks need large amounts of labeled training data to prevent overfitting; many real world IDS datasets lack sufficient numbers of labeled samples for certain rare types of attacks.
\item \textbf{Training Time Costs. }Large deep networks can take too long to train and/or make too expensive calls to perform real-time detection on commodity hardware.
\item \textbf{Black Box Interpretation. }In addition to the difficulty in understanding why a given flow is labeled as malicious, analysts require access to evidence of the attack to support investigations and demonstrate regulatory compliance.
\item \textbf{Fragility to Adversarial Attacks. }It is well established that deep networks are sensitive to adversarial examples---the inputs are slightly altered so as to result in incorrect classifications ~\cite{tafreshian2025adv}. 
\end{itemize}

\section{Hybrid and Ensemble Techniques. }


\subsection{Why do we want Hybrid Techniques? }

Because supervised methods are highly accurate at identifying previously seen attacks, but fail to recognize new attacks. Because unsupervised anomaly detection identifies new anomalies, but does so by labeling every new observation as potentially anomalous. A well-designed hybrid technique should utilize the strength of supervised classification for recognizing previously seen attacks and the sensitivity of unsupervised anomaly detection for recognizing new attacks.

\subsection{IDS Ensemble Methods. }

Khan et al.~\cite{khan2021invehicle} described an integration of multiple tier ensembles of Support Vector Machines (SVMs), Decision Trees and Logistic Regression through Stacking Meta-Learning. In this work, a two level classifier is trained to determine the best weights for aggregating the outputs of the base classifiers to minimize the total number of errors. Ring et al.~\cite{ring2019survey}, presented Dynamic Ensemble Selection (DES), a technique for selecting the best predictor from a pre-trained pool of predictors at test time based on the context provided by the test data. DES proved to be more resilient to changes in the input data distribution than traditional weighted average ensemble methods.

\subsection{Robustness Against Adversaries. }

Tafreshian and Zhang~\cite{tafreshian2025adv} discussed ways to defend against adversarial attacks on ML-based IDS systems; Tafreshian and Zhang note that as attackers begin to understand what methods are being employed against them, defending against such attacks will grow increasingly important. They propose a combination of adversarial training (where the original training dataset is augmented with adversarial examples) along with certified defense (guaranteeing that an attacker cannot perturb a network flow beyond a certain threshold). Sewak et al.~\cite{sewak2022drl}, also reviewed applying deep reinforcement learning (DRL) to cybersecurity; they demonstrated that DRL can enable continuous online adaptation in rapidly evolving cyber-threat environments where attacker tactics shift constantly due to defender actions.

\section{Positioning and Identified Research Gaps. }

From our review of the literature above we have identified three primary research gaps:

\begin{enumerate}
\item \textbf{Preprocessing Pipeline Integration. }Few researchers examine class imbalance, redundant features, and poor data quality jointly in a unified process; instead, researchers address these issues separately.
\item \textbf{Comparative Evaluation of Multiple Algorithms. }Most comparative evaluations include only a small number of algorithms; evaluation procedures tend to vary greatly between papers; few papers include complete reporting of meaningful metrics including both accuracy and f1 score as well as class-by-class analysis.
\item \textbf{Operational-Level Explanatory Power. }While SHAP has been shown applicable to various forms of cybersecurity tasks, the combined application of SHAP with LIME for both global and local interpretation in an IDS context with direct ties to analyst workflow processes has yet to be thoroughly documented.
\item \textbf{Deployment Architecture. }There is still a disconnect between evaluating algorithmic performance on benchmarks vs actual operational deployments; few authors describe viable architectures for deploying IDS solutions in commercial-grade SIEM systems.
\end{enumerate}

All of these gaps are addressed in the methodology portion of this research project.